%==============================================================================
% ICSE 2027 -- Tool Demonstration and Data Showcase Track
% arcade-agent: Architecture Recovery and Conformance Inside the Agent's Loop
%
% CFP facts verified 2026-10-07 against
%   https://conf.researchr.org/track/icse-2027/icse-2027-demonstrations
%   - 4 pages INCLUSIVE of references, figures, tables, appendices
%   - IEEEtran [10pt,conference]; single-anonymous (authors named)
%   - 3-5 min YouTube video, live through review; tool usable without building
%   - submission Fri 23 Oct 2026 AoE, https://icse27demos.hotcrp.com/
%
% Every number in Sec. V comes from eval/results/*.jsonl via eval/run.py
% --summarize; every number in Secs. III-IV from measure_budget.py.
%==============================================================================

\documentclass[10pt,conference]{IEEEtran}

\usepackage{cite}
\usepackage{graphicx}
\usepackage{booktabs}
\usepackage{xcolor}
\usepackage{listings}
\usepackage{amssymb}
\usepackage{balance}
\usepackage[hidelinks]{hyperref}
\usepackage{url}
\usepackage{tikz}
\usetikzlibrary{arrows.meta,positioning,fit,backgrounds}

\newif\ifdraftmode
\draftmodetrue

\ifdraftmode
  \newcommand{\todo}[1]{{\color{red}\small[\textbf{TODO} #1]}}
\else
  \newcommand{\todo}[1]{}
\fi

\newcommand{\tool}[1]{\texttt{#1}}
\newcommand{\arcade}{\textsc{arcade-agent}}

% ---- Evaluation numbers (filled from eval/results by run.py --summarize) ----
% Generated by make_results.py -- do not edit.
\newcommand{\HaikuOffViol}{20}
\newcommand{\HaikuOffN}{20}
\newcommand{\HaikuOffFunc}{20}
\newcommand{\HaikuOffUsed}{0}
\newcommand{\HaikuOffProposeCalls}{0}
\newcommand{\HaikuOffPreviewCalls}{0}
\newcommand{\HaikuOffCheckCalls}{0}
\newcommand{\HaikuOffCost}{0.06}
\newcommand{\HaikuOffWall}{38}
\newcommand{\SonnetOffViol}{1}
\newcommand{\SonnetOffN}{6}
\newcommand{\SonnetOffFunc}{6}
\newcommand{\SonnetOffUsed}{0}
\newcommand{\SonnetOffProposeCalls}{0}
\newcommand{\SonnetOffPreviewCalls}{0}
\newcommand{\SonnetOffCheckCalls}{0}
\newcommand{\SonnetOffCost}{0.05}
\newcommand{\SonnetOffWall}{20}
\newcommand{\HaikuSpecViol}{5}
\newcommand{\HaikuSpecN}{20}
\newcommand{\HaikuSpecFunc}{20}
\newcommand{\HaikuSpecUsed}{0}
\newcommand{\HaikuSpecProposeCalls}{0}
\newcommand{\HaikuSpecPreviewCalls}{0}
\newcommand{\HaikuSpecCheckCalls}{0}
\newcommand{\HaikuSpecCost}{0.09}
\newcommand{\HaikuSpecWall}{59}
\newcommand{\SonnetSpecViol}{0}
\newcommand{\SonnetSpecN}{6}
\newcommand{\SonnetSpecFunc}{6}
\newcommand{\SonnetSpecUsed}{0}
\newcommand{\SonnetSpecProposeCalls}{0}
\newcommand{\SonnetSpecPreviewCalls}{0}
\newcommand{\SonnetSpecCheckCalls}{0}
\newcommand{\SonnetSpecCost}{0.07}
\newcommand{\SonnetSpecWall}{23}
\newcommand{\HaikuAdvisoryViol}{0}
\newcommand{\HaikuAdvisoryN}{20}
\newcommand{\HaikuAdvisoryFunc}{20}
\newcommand{\HaikuAdvisoryUsed}{0}
\newcommand{\HaikuAdvisoryProposeCalls}{0}
\newcommand{\HaikuAdvisoryPreviewCalls}{0}
\newcommand{\HaikuAdvisoryCheckCalls}{0}
\newcommand{\HaikuAdvisoryCost}{0.11}
\newcommand{\HaikuAdvisoryWall}{68}
\newcommand{\SonnetAdvisoryViol}{0}
\newcommand{\SonnetAdvisoryN}{6}
\newcommand{\SonnetAdvisoryFunc}{6}
\newcommand{\SonnetAdvisoryUsed}{0}
\newcommand{\SonnetAdvisoryProposeCalls}{0}
\newcommand{\SonnetAdvisoryPreviewCalls}{0}
\newcommand{\SonnetAdvisoryCheckCalls}{0}
\newcommand{\SonnetAdvisoryCost}{0.06}
\newcommand{\SonnetAdvisoryWall}{18}
\newcommand{\HaikuToolsViol}{0}
\newcommand{\HaikuToolsN}{20}
\newcommand{\HaikuToolsFunc}{20}
\newcommand{\HaikuToolsUsed}{20}
\newcommand{\HaikuToolsProposeCalls}{81}
\newcommand{\HaikuToolsPreviewCalls}{64}
\newcommand{\HaikuToolsCheckCalls}{40}
\newcommand{\HaikuToolsCost}{0.15}
\newcommand{\HaikuToolsWall}{88}
\newcommand{\SonnetToolsViol}{0}
\newcommand{\SonnetToolsN}{6}
\newcommand{\SonnetToolsFunc}{6}
\newcommand{\SonnetToolsUsed}{6}
\newcommand{\SonnetToolsProposeCalls}{13}
\newcommand{\SonnetToolsPreviewCalls}{8}
\newcommand{\SonnetToolsCheckCalls}{5}
\newcommand{\SonnetToolsCost}{0.07}
\newcommand{\SonnetToolsWall}{24}
\newcommand{\HaikuOnViol}{0}
\newcommand{\HaikuOnN}{20}
\newcommand{\HaikuOnFunc}{20}
\newcommand{\HaikuOnUsed}{20}
\newcommand{\HaikuOnProposeCalls}{84}
\newcommand{\HaikuOnPreviewCalls}{59}
\newcommand{\HaikuOnCheckCalls}{36}
\newcommand{\HaikuOnCost}{0.14}
\newcommand{\HaikuOnWall}{79}
\newcommand{\SonnetOnViol}{0}
\newcommand{\SonnetOnN}{6}
\newcommand{\SonnetOnFunc}{6}
\newcommand{\SonnetOnUsed}{6}
\newcommand{\SonnetOnProposeCalls}{13}
\newcommand{\SonnetOnPreviewCalls}{7}
\newcommand{\SonnetOnCheckCalls}{6}
\newcommand{\SonnetOnCost}{0.05}
\newcommand{\SonnetOnWall}{26}
\newcommand{\EvalTable}{%
\begin{tabular}{@{}lrrrr@{}}
\toprule
Condition & \multicolumn{2}{c}{Haiku 4.5} & \multicolumn{2}{c}{Sonnet 5.5} \\
 & viol. & func. & viol. & func. \\
\midrule
\emph{off} (no spec) & 20/20 & 20/20 & 1/6 & 6/6 \\
spec & 5/20 & 20/20 & 0/6 & 6/6 \\
spec + rules & 0/20 & 20/20 & 0/6 & 6/6 \\
spec + tools & 0/20 & 20/20 & 0/6 & 6/6 \\
spec + rules + tools & 0/20 & 20/20 & 0/6 & 6/6 \\
\bottomrule
\end{tabular}}
\newcommand{\ModelNames}{Claude Haiku 4.5 and Claude Sonnet 5.5}
\newcommand{\NumRuns}{130}
\newcommand{\OffViol}{20}
\newcommand{\OffN}{20}
\newcommand{\PSpecVsTools}{0.024}
\newcommand{\PSpecVsAdvisory}{0.024}
\newcommand{\PSpecVsOn}{0.024}
\newcommand{\InjectRules}{api imports store: forbidden-dependency; store imports service: dependency-cycle, layer-violation; domain imports store: dependency-cycle, forbidden-dependency}
\newcommand{\InjectCaught}{3}
\newcommand{\InjectN}{3}
\newcommand{\InjectBaseline}{PASS}
\newcommand{\BudgetTable}{\todo{run measure\_budget.py}}
\newcommand{\SelfFullTok}{\todo{n}}
\newcommand{\SelfSummTok}{\todo{n}}
\newcommand{\SelfRawTok}{\todo{n}}
\newcommand{\CtxFiles}{\todo{n}}
\newcommand{\CtxTok}{\todo{n}}
\newcommand{\TokenFootnote}{\todo{tokenizer footnote}}
\newcommand{\ResultSentence}{\todo{ResultSentence}}
\newcommand{\EvidenceProse}{\todo{EvidenceProse}}
\newcommand{\TranscriptListing}{\todo{TranscriptListing}}


\lstset{
  basicstyle=\ttfamily\scriptsize,
  breaklines=true,
  columns=fullflexible,
  frame=single,
  framesep=3pt,
  aboveskip=4pt,
  belowskip=4pt,
  xleftmargin=2pt,
  literate={->}{{$\rightarrow$}}1
}

\tikzset{
  comp/.style={draw, rounded corners=2pt, minimum width=1.25cm, minimum height=0.5cm,
               font=\scriptsize\ttfamily, fill=white},
  dep/.style={-{Stealth[length=4pt]}, semithick},
  bad/.style={dep, red!75!black, very thick},
  box/.style={draw, rounded corners=2pt, align=center, font=\scriptsize,
              minimum height=0.55cm, fill=white},
}

\begin{document}

\title{\arcade{}: Architecture Recovery and Conformance\\Inside the Coding Agent's Loop}

\author{
  \IEEEauthorblockN{Duc Minh Le\todo{affiliation, email}}
  \and
  \IEEEauthorblockN{\todo{Tuan's full name}}
  \IEEEauthorblockA{\todo{affiliation, email}}
}

\maketitle

%==============================================================================
\begin{abstract}
Coding agents now write much of the code in many projects, and nothing in the
agent's loop guards its architecture: type checkers, linters and tests run
continuously, while architecture conformance is checked in CI, if at all, after
the code exists. We present \arcade{}, an installable library and Model Context
Protocol server that makes architecture recovery, smell detection, decay
metrics and conformance checking affordable for an agent to call while it
works. Session handles and per-call token budgets keep responses small, and a
deterministic guardrail answers \emph{where should this code go} and \emph{may
it depend on that} before the code is written. In a controlled experiment on a
greenfield task, \ResultSentence{} The tool is MIT-licensed and installs with
\tool{pip install arcade-agent}.

\smallskip
\noindent Video: \url{\todo{YouTube URL}}
\end{abstract}

\begin{IEEEkeywords}
software architecture, architecture recovery, conformance checking, AI coding
agents, Model Context Protocol
\end{IEEEkeywords}

%==============================================================================
\section{Introduction}
%==============================================================================
Software engineering rests on an assumption so basic that it is rarely written
down: the artifacts under discussion are ultimately produced by people. The
discipline's vocabulary inherits that premise everywhere. Effort is denominated
in the man-month~\cite{brooks1975mythical}; productivity was once counted in
lines written per day. The unit has changed and the premise has not --- teams now
compare the tokens their agents consume, a measure that captures activity rather
than progress in exactly the way its predecessor did.

Coding agents have broken the assumption. An agent no longer completes a line;
it plans a change, edits across files, runs the tests, and
iterates~\cite{sweagent,swebench}. What was built on the premise that a person
stands at every step --- the process, the review points, the tools, and the
guarantees each was trusted to provide --- is now due for re-examination.

This paper takes up one instance: what keeps a system's structure coherent while
that system is being built. Type checkers guard types, linters guard style, test
suites guard behavior. All three sit inside the agent's loop, and an agent
consults them continuously and cheaply. Nothing plays that role for
architecture. Conformance checking is a mature
idea~\cite{murphy1995reflexion,passos2010conformance,archunit}, but it runs in
continuous integration, after the code is written, addressed to a human
reviewer. For an agent that is both too late and too blunt: the code already
exists, and a failed check that carries no remedy produces thrashing rather than
repair. Meanwhile the decay such a check is meant to catch is measurable and
predicts where defects and costly change will
land~\cite{le2018decay,le2021decay-predictor}.

We present \arcade{}, which moves architecture recovery, smell detection, decay
metrics, and conformance checking inside the agent's loop. The analyses are
established --- most come from the ARCADE line of work~\cite{laser2020arcade} ---
and we claim no novelty for them. Our contribution is making them something an
agent can afford to call, and can call \emph{before} the code exists:
\begin{itemize}
  \item an installable, framework-agnostic analysis library with seven language
        parsers and five recovery algorithms;
  \item a Model Context Protocol (MCP)~\cite{mcp} server whose 24 tools return session handles and respect a
        per-call token budget, so an agent receives summaries, not dumps;
  \item task-shaped context tools that answer an agent's actual questions ---
        what to read, what a change breaks --- instead of returning a graph;
  \item a deterministic architecture guardrail, available as MCP tools and as a
        pre-commit/CI gate, and a controlled experiment that separates what it
        adds from simply stating the rules.
\end{itemize}

%==============================================================================
\section{Motivating Scenario}
\label{sec:motivation}
%==============================================================================
An agent is asked to add order placement to a new Python service whose folders
already name the layers: \tool{api}, \tool{service}, \tool{store} and
\tool{domain}. The team's rule is that the API never touches the store directly.
The task also says that \tool{OrdersApi()} must work without arguments, so the
API class has to wire up its own collaborators.

The shortest path to a working feature is for \tool{OrdersApi} to construct the
store itself, or to skip the service layer entirely. Both pass every test. In
our experiment (Sec.~\ref{sec:evidence}) an agent without architectural guidance
took one of these paths in \OffViol{} of \OffN{} runs (Fig.~\ref{fig:scenario},
left). With the guardrail it calls \tool{propose\_placement} before writing,
learns that persistence belongs in \tool{store} and that \tool{api} may depend
only on \tool{service}, checks the edge it is about to add with
\tool{preview\_impact}, and routes through the service instead
(Fig.~\ref{fig:scenario}, right).

\begin{figure}[t]
  \centering
  \begin{tikzpicture}[node distance=0.42cm and 0.5cm]
    \begin{scope}
      \node[comp] (a1) {api};
      \node[comp, below=of a1] (s1) {service};
      \node[comp, below left=0.42cm and -0.25cm of s1] (st1) {store};
      \node[comp, below right=0.42cm and -0.25cm of s1] (d1) {domain};
      \draw[dep] (a1) -- (s1);
      \draw[bad] (a1.west) to[bend right=35] node[left, font=\tiny, text=red!75!black]{forbidden} (st1.west);
      \draw[dep] (s1) -- (st1);
      \draw[dep] (st1) -- (d1);
      \node[above=0.08cm of a1, font=\scriptsize\itshape] {unaided};
    \end{scope}
    \begin{scope}[xshift=4.1cm]
      \node[comp] (a2) {api};
      \node[comp, below=of a2] (s2) {service};
      \node[comp, below left=0.42cm and -0.25cm of s2] (st2) {store};
      \node[comp, below right=0.42cm and -0.25cm of s2] (d2) {domain};
      \draw[dep] (a2) -- (s2);
      \draw[dep] (s2) -- (st2);
      \draw[dep] (s2) -- (d2);
      \draw[dep] (st2) -- (d2);
      \node[above=0.08cm of a2, font=\scriptsize\itshape] {with the guardrail};
    \end{scope}
  \end{tikzpicture}
  \caption{Component dependencies an agent produced for the same task. Left:
  the API constructs the store itself. Right: the API goes through the
  service, as the specification requires.}
  \label{fig:scenario}
\end{figure}

%==============================================================================
\section{Architecture and Implementation}
\label{sec:impl}
%==============================================================================
\arcade{} is a library first and a server second. The analysis core is a
sequence of pure stages --- \tool{ingest}, \tool{parse}, \tool{recover},
\tool{detect\_smells}, \tool{compute\_metrics}, \tool{compare} --- each
independently callable and each returning a plain dataclass. The MCP server,
the guardrail and the CI entry points consume that core rather than wrap it
(Fig.~\ref{fig:pipeline}).

\begin{figure}[t]
  \centering
  \begin{tikzpicture}[node distance=0.3cm]
    \node[box, minimum width=7.9cm] (agent) {coding agent (Claude Code, Cursor, \dots)};
    \node[box, below=0.35cm of agent, xshift=-2.05cm, minimum width=3.8cm] (mcp)
      {MCP server \tool{arcade-mcp}\\[-1pt]\tiny session handles $\cdot$ token budgets};
    \node[box, below=0.35cm of agent, xshift=2.05cm, minimum width=3.8cm] (cli)
      {\tool{arcade-guard} CLI\\[-1pt]\tiny hook $\cdot$ pre-commit $\cdot$ CI gate};
    \node[box, below=0.35cm of mcp, minimum width=3.8cm] (an)
      {analysis tools\\[-1pt]\tiny summarize, context\_for\_task, diff\_impact, \dots};
    \node[box, below=0.35cm of cli, minimum width=3.8cm] (gd)
      {conformance engine\\[-1pt]\tiny propose $\cdot$ preview $\cdot$ check};
    \node[box, below=0.35cm of an, xshift=2.05cm, minimum width=7.9cm] (core)
      {core: 7 parsers $\rightarrow$ dependency graph $\rightarrow$ recovery (pkg, wca, acdc, arc, limbo) $\rightarrow$ smells, metrics};
    \node[box, right=0.15cm of gd.north east, anchor=north west, minimum width=1.1cm,
          fill=yellow!12] (spec) {\tiny architecture\\[-2pt]\tiny\tool{.spec.json}};
    \draw[dep] (agent) -- (mcp);
    \draw[dep] (agent) -- (cli);
    \draw[dep] (mcp) -- (an);
    \draw[dep] (mcp.east) -- (gd.north west);
    \draw[dep] (cli) -- (gd);
    \draw[dep] (an) -- (an |- core.north);
    \draw[dep] (gd) -- (gd |- core.north);
    \draw[dep, dashed] (spec) -- (gd);
  \end{tikzpicture}
  \caption{\arcade{} in the agent's loop. The MCP tools and the CLI gate call
  one conformance engine, so their verdicts cannot disagree.}
  \label{fig:pipeline}
\end{figure}

\subsection{Parsing and recovery}
Each language is a \tool{LanguageParser} registered by decorator. Seven are
available: Java and Python in the core install; C, TypeScript, Go, Kotlin and
Rust with the \tool{[languages]} extra. Parsing is two-pass --- a per-file AST
walk, then a link pass that resolves imports, inheritance and references across
the file set. Polyglot repositories are merged into one graph, but cross-language
edges are linked only within a language family; Java and Kotlin are the one
family we have validated. Five recovery algorithms are available ---
package-based grouping (\tool{pkg}), \tool{wca}, \tool{acdc}, and the
LLM-assisted \tool{arc} and \tool{limbo}. They are simplified reimplementations
of prior
work~\cite{tzerpos2000acdc,garcia2011arc,andritsos2004limbo,maqbool2004wca};
\tool{arc} and \tool{limbo} replace topic models with LLM-assigned concern
labels and are nondeterministic.

\subsection{Delivering analysis an agent can afford}
Architecture analysis produces far more data than an agent can read: the full
dependency graph and recovered architecture of \arcade{} itself costs
\textbf{\SelfFullTok{} tokens}.\footnote{\TokenFootnote} Two mechanisms make it
affordable. First, every tool returns a compact summary with a
\tool{session\_id}; the analysed object stays server-side, downstream tools
accept the ID, and \tool{get\_full\_result} retrieves it only on request, so an
agent pipes \tool{parse} into \tool{recover} into \tool{detect\_smells} while
the graph never crosses the wire. Second, every tool accepts \tool{max\_tokens}.
Payloads are reduced by ordered strategies --- truncate component member lists,
collapse entities to an FQN$\rightarrow$kind map, collapse edges to per-relation
counts, drop maps for counts --- re-checking the budget after each, so the
cheapest sufficient reduction wins and the response is flagged as partial.
Table~\ref{tab:budget} shows the effect.

\begin{table}[t]
  \centering
  \caption{Tokens to load an architecture, by repository and access path.}
  \label{tab:budget}
  \footnotesize
  \BudgetTable
\end{table}

\subsection{A deterministic guardrail}
The guardrail checks code against an author-written
\tool{architecture.spec.json}: components defined by path globs and assigned to
layers, allowed and forbidden dependencies (each with a reason), and budgets
such as no cycles or a fan-in ceiling. Entities are mapped to components by
path, not by clustering, and rules are evaluated over the aggregated
component edges, so the same code always receives the same verdict. A guardrail
an agent consults dozens of times per session must be deterministic: an agent
that gets a different answer on retry thrashes. One function decides whether a
component edge is allowed; \tool{preview\_impact} applies it to a hypothetical
edge and \tool{check\_architecture} to every actual one, so the answer before
the import and the verdict after it cannot disagree. The same engine runs as
\tool{arcade-guard check}, which exits non-zero on failure for a pre-commit hook
or CI job, and as an agent hook after each edit. Conformance is computed over
\emph{referenced} dependencies: an import that is never used, or used only in a
type annotation, does not create an edge.

%==============================================================================
\section{The Tool in Action}
\label{sec:action}
%==============================================================================
The demonstration video follows one agent session end to end.

\paragraph*{Setup}
\tool{pip install "arcade-agent[mcp,languages]"} installs the library, the
\tool{arcade-mcp} server, which the agent launches over stdio, and the
\tool{arcade-guard} CLI. Reading a codebase needs no configuration; the
guardrail needs a spec, which \tool{init\_spec} scaffolds from a layered,
hexagonal, clean or MVC template.

\paragraph*{Understand}
An agent's first question about an unfamiliar repository is what is in it.
\tool{summarize} answers with the package tree, dependency hotspots, entry
points and entity-kind counts in one call --- \SelfSummTok{} tokens on \arcade{}
itself, against \SelfRawTok{} to read its source.

\paragraph*{Target}
Next the agent needs to know what to read. Given the task \emph{``add a Ruby
language parser''}, \tool{context\_for\_task} returns \CtxFiles{} ranked files in
\CtxTok{} tokens, each with a role and a reason; the highest-ranked are the
existing parsers a contributor would open first. Ranking is lexical, over entity
names, packages and paths, boosted by component membership, so it degrades when
the task's vocabulary and the code's diverge.

\paragraph*{Guard}
Before writing, the agent calls \tool{propose\_placement} with a description of
what it is about to add and learns the component it belongs in and what that
component may depend on. Before a cross-component import it calls
\tool{preview\_impact}. Fig.~\ref{fig:transcript} shows a session from our
experiment in which the agent asks before it writes.

\paragraph*{Gate}
After the edit, \tool{check\_architecture} returns PASS, WARN or FAIL, each
violation carrying the spec's reason and a fix. The same check gates commits and
pull requests. \arcade{} also analyses itself on every push and commits the
recovered architecture as a baseline, so each pull request is compared with the
last accepted state.

\begin{figure}[t]
\TranscriptListing
  \caption{Excerpt of a recorded agent session (\tool{tools} condition,
  Sec.~\ref{sec:evidence}); agent commands and guardrail output, abridged.}
  \label{fig:transcript}
\end{figure}

%==============================================================================
\section{Evidence It Works}
\label{sec:evidence}
%==============================================================================
Does the guardrail change what an agent builds, and is that more than the
effect of telling it the rules? We ran Claude Code headlessly on the greenfield
task of Sec.~\ref{sec:motivation} --- an empty \tool{app/\{api, service, store,
domain\}} scaffold, with folder roles stated in every prompt --- under five
conditions. \emph{off} has no specification. The other four place the spec in
the repository and vary two factors: whether the prompt states the dependency
rules, and whether the agent is given the guardrail workflow (\tool{propose},
\tool{preview}, then \tool{check} until PASS). User settings, plugins, hooks and
MCP servers were disabled, so only the prompt differed. Each result was scored
by \tool{check\_architecture} against the canonical spec, counting only
error-severity violations, and by a functional test of the feature.
Table~\ref{tab:eval} reports the results.

\begin{table}[t]
  \centering
  \caption{Agent runs violating the specification (errors only) and passing
  the functional test, by condition.}
  \label{tab:eval}
  \footnotesize
  \EvalTable
\end{table}

\EvidenceProse

\paragraph*{Detection}
Separately, we injected three shortcuts into a conformant version of the
service --- the API importing the store, the store importing the service, the
domain importing the store. \tool{check\_architecture} reported each with the
expected rule (\InjectRules{}) and \tool{arcade-guard check} exited 1 every
time.

\paragraph*{Threats to validity}
One task, one language, and one agent harness: the result shows the mechanism
works, not how often it matters in practice. Folder roles were stated in every
condition, which nudges toward layering, so the unaided rate is likely a floor.
The models are \ModelNames{}; behaviour differs across models, and Sonnet's
lower violation rate leaves less room for the guardrail to help. Conformance is
structural; the functional test checks the feature, not code quality.

%==============================================================================
\section{Related Work}
\label{sec:related}
%==============================================================================
\textbf{Architecture recovery and decay.} ACDC, ARC, LIMBO and WCA cluster code
into components~\cite{tzerpos2000acdc,garcia2011arc,andritsos2004limbo,maqbool2004wca};
ARCADE packaged them with decay metrics for
researchers~\cite{laser2020arcade,behnamghader2017evolution,le2015change}.
\arcade{} reimplements a subset for a different consumer --- an agent, mid-task,
with a token budget.
\textbf{Conformance checking.} Reflexion models compare a source model with an
intended one~\cite{murphy1995reflexion}, and tools such as
ArchUnit~\cite{archunit} enforce rules in tests~\cite{passos2010conformance}.
These are invoked by people, after the code exists; our guardrail is consulted
by the agent before and while it writes, and shares one engine with the gate.
Delivering analysis at the moment of change is what made static analysis
effective at scale~\cite{tricorder}; we apply that lesson above the code level.
\textbf{Context for coding agents.} Repository graphs and code-graph queries
improve agents' task success~\cite{repograph,codexgraph}, and simple pipelines
rival agents~\cite{agentless}; these optimize \emph{what to read} and measure
correctness. Recent work reports that agent-written code accumulates
architectural smells~\cite{zhu2026decay}, and that agents asked to repair such smells introduce new ones~\cite{smellbench2026}. \arcade{} adds an architectural
contract the agent can query and measures conformance of its output.

%==============================================================================
\section{Availability and Limitations}
\label{sec:conclusion}
%==============================================================================
\arcade{} is MIT-licensed:
\tool{pip install "arcade-agent[mcp,languages]"} (Python $\ge$3.12). Source,
documentation and the evaluation harness with all run records are at
\url{https://github.com/arcade-agent/arcade-agent} and
\url{https://arcade-agent.dev}. Limitations: cross-language linking is
validated only for Java and Kotlin; parsers emit import, inheritance and
reference edges, not call edges; incremental parsing covers Python only; and the
guardrail checks structure, not behaviour.

\balance
\bibliographystyle{IEEEtran}
\bibliography{refs}

\end{document}
